% Options for packages loaded elsewhere
\PassOptionsToPackage{unicode}{hyperref}
\PassOptionsToPackage{hyphens}{url}
\documentclass[
  11pt,
]{article}
\usepackage{xcolor}
\usepackage[margin=25mm]{geometry}
\usepackage{amsmath,amssymb}
\setcounter{secnumdepth}{-\maxdimen} % remove section numbering
\usepackage{iftex}
\ifPDFTeX
  \usepackage[T1]{fontenc}
  \usepackage[utf8]{inputenc}
  \usepackage{textcomp} % provide euro and other symbols
\else % if luatex or xetex
  \usepackage{unicode-math} % this also loads fontspec
  \defaultfontfeatures{Scale=MatchLowercase}
  \defaultfontfeatures[\rmfamily]{Ligatures=TeX,Scale=1}
\fi
\usepackage{lmodern}
\ifPDFTeX\else
  % xetex/luatex font selection
\fi
% Use upquote if available, for straight quotes in verbatim environments
\IfFileExists{upquote.sty}{\usepackage{upquote}}{}
\IfFileExists{microtype.sty}{% use microtype if available
  \usepackage[]{microtype}
  \UseMicrotypeSet[protrusion]{basicmath} % disable protrusion for tt fonts
}{}
\makeatletter
\@ifundefined{KOMAClassName}{% if non-KOMA class
  \IfFileExists{parskip.sty}{%
    \usepackage{parskip}
  }{% else
    \setlength{\parindent}{0pt}
    \setlength{\parskip}{6pt plus 2pt minus 1pt}}
}{% if KOMA class
  \KOMAoptions{parskip=half}}
\makeatother
\setlength{\emergencystretch}{3em} % prevent overfull lines
\providecommand{\tightlist}{%
  \setlength{\itemsep}{0pt}\setlength{\parskip}{0pt}}
\usepackage{bookmark}
\IfFileExists{xurl.sty}{\usepackage{xurl}}{} % add URL line breaks if available
\urlstyle{same}
\hypersetup{
  hidelinks,
  pdfcreator={LaTeX via pandoc}}

\author{}
\date{}

\begin{document}

\section{Kasparov's technical
theorem}\label{kasparovs-technical-theorem}

\emph{Written by GPT-6.1 Sol (OpenAI). Public domain (CC0).}

Constructing a Kasparov product requires combining two operators whose
errors are compact in different directions. A partition of the identity
in a multiplier algebra lets us assign one error to one part and the
other error to the complementary part. The partition must also commute
modulo the ideal with the operators already present. Kasparov's
technical theorem supplies exactly this partition.

We prove the theorem, including its graded form, by constructing a
strict sum of positive operators. The proof includes the required
convexity argument, the reduction from \(\sigma\)-unital to separable
data, and the estimates for each error. We also prove the multiplier
quotient lifting used in a shorter separation argument when no
commutator control is requested. Finally we use the theorem to extend an
operator homotopy from an ideal; this is the technical input to the
comparison of cycle relations.

We use the earlier continuous functional-calculus proofs. Theorem 1.1
proves its bounded-functional and strong-support inputs locally, using
the GNS construction of
\href{KT-KK-05.html\#lemma-3-0a-the-ordinary-spatial-norm}{Lesson 05,
Lemma 3.0a}. The ideal version of quasicentral approximate identities
belongs to \emph{Lifting completely positive maps}; Section 1 supplies
the additional derivating-subspace version needed here. For a countably
generated Hilbert module, the identification
\(\mathcal L(E)=M(\mathcal K(E))\) and \(\sigma\)-unitality of
\(\mathcal K(E)\) are proved in
\href{KT-KK-05.html\#lemma-2-0-compact-module-operators-and-their-multipliers}{Lesson
05, Lemma 2.0}.

\subsection{1. Approximate identities with commutator
control}\label{approximate-identities-with-commutator-control}

An approximate identity means positive contractions \(u_\lambda\) such
that \(u_\lambda a\to a\) and \(au_\lambda\to a\) in norm for every
\(a\) in the algebra. A subspace \(\Delta\) of an ambient C*-algebra
\textbf{derives} a subalgebra \(A\) if \[
[d,a]\in A\qquad(d\in\Delta,\ a\in A).
\tag{1.1}
\] Here the commutator is ordinary until grading is explicitly
introduced. Deriving is weaker than preserving \(A\) by left and right
multiplication separately.

\textbf{Theorem 1.1 (convex quasicentrality).} Suppose \(\Delta\)
derives \(A\). From any positive contractive approximate identity of
\(A\), finite convex combinations can be chosen to form an approximate
identity quasicentral for \(\Delta\): \[
\|[u_\lambda,d]\|\longrightarrow0\qquad(d\in\Delta).
\tag{1.2}
\] More precisely every tail of the original approximate identity has
convex combinations meeting any finite collection of
approximate-identity and commutator tests. If \(A\) is \(\sigma\)-unital
and \(\Delta\) is separable, the combinations may be chosen as a
sequence.

\textbf{Proof.} Write \(C\) for the ambient C*-algebra. First every
bounded complex linear functional \(f\) on \(C\) is a vector coefficient
in some Hilbert-space representation of the forced unitization \(C^+\).
Here is a proof that needs only Hahn--Banach and the GNS construction
proved in Lesson 05 Lemma 3.0a. The zero functional is immediate;
normalize a nonzero one to \(\|f\|=1\). In \(M_2(C^+)\) use the unital
self-adjoint subspace consisting of matrices \[
\begin{pmatrix}\lambda1&a\\b&\mu1\end{pmatrix},
\qquad a,b\in C,\quad\lambda,\mu\in\mathbb C.
\] On it define \[
F\left(\begin{pmatrix}\lambda1&a\\b&\mu1\end{pmatrix}\right)
=\tfrac12(\lambda+\mu+f(a)+\overline{f(b^*)}).
\] This is complex linear and real on self-adjoint matrices. For a
positive such matrix, \(\lambda,\mu\geq0\) and
\(\|a\|^2\leq\lambda\mu\). To check that bound when \(\lambda>0\), apply
the triangular Schur-complement factorization to obtain
\(a^*a\leq\lambda\mu1\). If \(\lambda=0\), positivity and its square
root show that the first column is zero, so \(a=0\). Thus \[
F\left(\begin{pmatrix}\lambda1&a\\a^*&\mu1\end{pmatrix}\right)
\geq(\lambda+\mu)/2-\|a\|\geq0.
\] Since \(F(1)=1\), order bounds give real norm one on the self-adjoint
part of this subspace. Extend it by real Hahn--Banach to a norm-one real
functional on \(M_2(C^+)_{\rm sa}\), retaining its value one at the
identity. This extension is positive: for \(0\leq h\leq1\), the bound
\(\|1-h\|\leq1\) gives \(F(h)=1-F(1-h)\geq0\), and scale for arbitrary
positive \(h\). Complexification therefore gives a state \(\Psi\) on
\(M_2(C^+)\).

Let \((\rho,H,\xi)\) be its cyclic GNS representation. Restrict \(\rho\)
to the diagonal copy of \(C^+\), writing
\(\pi(c)=\rho(\operatorname{diag}(c,c))\). If \(E_{ij}\) are the scalar
matrix units, the prescribed off-diagonal value is \[
f(c)/2=\Psi(E_{12}\operatorname{diag}(c,c))
=\langle\rho(E_{21})\xi,\pi(c)\xi\rangle\qquad(c\in C).
\] Rescale the first vector for the original norm of \(f\). This
establishes the asserted vector-coefficient representation; it uses
neither a polar decomposition of functionals nor a Jordan decomposition.

For this representation \(\pi\), let \(H_A=\overline{\pi(A)H}\), and let
\(P\) be its Hilbert-space orthogonal projection. The projection needs
only the Hilbert-space projection theorem: minimize distance to a closed
linear subspace, use the parallelogram identity to make a minimizing
sequence Cauchy, and differentiate the squared distance along each
subspace direction to get orthogonality. This also proves the resulting
decomposition and bounded projection. Approximate-identity convergence
shows \(\pi(e_\lambda)\to P\) strongly: check on \(\pi(a)\eta\) by norm
convergence \(e_\lambda a\to a\), extend by their dense span and the
contraction bound, and note that each \(\pi(e_\lambda)\) vanishes on
\(H_A^\perp\).

The derivating hypothesis also holds for \(d^*\), since
\([d^*,a]=-[d,a^*]^*\). Both \(\pi(d)\) and \(\pi(d^*)\) preserve
\(H_A\), because \[
\pi(d)\pi(a)\eta
=\pi(a)\pi(d)\eta+\pi([d,a])\eta\in H_A.
\] Hence \(H_A\) reduces \(\pi(d)\) and \([P,\pi(d)]=0\). Strong
convergence now gives \(\pi([e_\lambda,d])\to0\) strongly. Testing its
vector coefficient proves \(f([e_\lambda,d])\to0\) for every bounded
functional \(f\), which is the claimed weak convergence in \(C\).

Fix \(d_1,\ldots,d_r\). The tuples
\(([u_\lambda,d_1],\ldots,[u_\lambda,d_r])\) converge weakly to zero in
the finite direct sum of the ambient algebra. Zero consequently lies in
the norm-closed convex hull of every tail. Indeed, otherwise
Hahn--Banach separation would give a bounded linear functional strictly
separating zero from that convex hull, contradicting weak convergence. A
finite convex combination in the tail therefore makes every indicated
commutator smaller than any prescribed positive tolerance.

The approximate-identity and commutator requirements can be imposed in
one separation step. For finite \(a_1,\ldots,a_s\in A\), adjoin the
components \(u_\lambda a_j-a_j\) and \(a_j u_\lambda-a_j\) to that
commutator tuple. These new components tend to zero in norm, while the
commutator components tend weakly to zero. The enlarged tuple therefore
has weak limit zero, and zero belongs to the norm-closed convex hull of
every tail of these enlarged tuples. A single finite convex combination
makes every component smaller than the chosen tolerance. Its associated
combination of the \(u_\lambda\) remains a positive contraction, and
linearity gives exactly the tested approximate-identity and commutator
errors. Directing these choices by finite tests, tails and tolerances
proves the net assertion. For a sequence, use a countable approximate
identity, a countable dense family in \(\Delta\), and progressively all
previous tests. The uniform bounds \(\|u_\lambda\|\leq1\) and
\(\|[u_\lambda,d]\|\leq2\|d\|\) extend convergence from the dense family
to all \(d\). \(\square\)

The same argument works uniformly on a norm-compact set of tests, by
choosing a finite net in that set.

We need increasing choices whose square-root differences have small
commutators. A \(\sigma\)-unital algebra has a strictly positive element
\(h\), equivalently one for which \(hA\) is norm dense. Choose
continuous functions \(f_n\) on \([0,\|h\|]\), vanishing near zero, with
\(0\leq f_n\leq1\), such that \[
e_n=f_n(h),\quad e_{n+1}e_n=e_n,\quad e_n\longrightarrow1
\quad\hbox{strictly}.
\tag{1.3}
\] For example let \(f_n\) vanish below \(2^{-n}\), equal one above
\(2^{-n+1}\), and increase linearly between, after rescaling \(h\). The
support condition gives the product equation. Functional calculus and
density of \(hA\) prove the approximate-identity assertion.

\textbf{Lemma 1.2 (increasing quasicentral choices).} Finite convex
combinations of the \(e_n\) can be chosen as an approximate identity
\(v_n\) with \[
v_nv_{n-1}=v_{n-1},\qquad 0\leq v_{n-1}\leq v_n\leq1,
\tag{1.4}
\] while satisfying successively prescribed finite or compact commutator
tests with a separable derivating subspace. The same assertion holds for
commutator tests on the images of \(v_n\) in a quotient algebra.

\textbf{Proof.} Having chosen \(v_{n-1}\) as a combination of finitely
many \(e_j\), choose all indices in the next combination greater than
their maximum. By (1.3), each new \(e_j\) multiplies every old \(e_i\)
to \(e_i\) on both sides. Hence the product equation in (1.4) holds. For
commuting positive contractions it implies \(v_n\geq v_{n-1}\). Theorem
1.1 provides the desired new finite tests from this tail. Let the
minimum index tend to infinity to retain the approximate-identity
property.

For quotient tests apply Theorem 1.1 to the image approximate identity,
selecting its convex coefficients. Use those same coefficients on the
original \(e_j\). Their quotient then has exactly the chosen commutator
bounds, while the original product relation and tail bounds are
retained. \(\square\)

\textbf{Lemma 1.3 (square-root commutators).} Given \(\varepsilon>0\),
there is \(\eta>0\) such that, in any C*-algebra, \[
0\leq a\leq1,\quad\|w\|\leq1,\quad
\|[a,w]\|<\eta
\quad\Longrightarrow\quad
\|[a^{1/2},w]\|<\varepsilon.
\tag{1.5}
\]

\textbf{Proof.} Approximate \(t^{1/2}\) uniformly on \([0,1]\) by a
polynomial \(p(t)=\sum c_kt^k\), to error less than \(\varepsilon/4\).
The commutator difference between \(a^{1/2}\) and \(p(a)\) has norm at
most \(\varepsilon/2\). Also \[
[a^k,w]=\sum_{j=0}^{k-1}a^j[a,w]a^{k-1-j},
\] so \(\|[p(a),w]\|\leq(\sum k|c_k|)\|[a,w]\|\). Taking \(\eta\) small
proves (1.5), uniformly in the algebra. \(\square\)

\subsection{2. The separation theorem}\label{the-separation-theorem}

The \textbf{strict topology} of \(M(J)\) tests convergence after
multiplication by each \(z\in J\), on both sides. A bounded strict
Cauchy sequence has a strict limit: define its left and right actions on
\(J\) by the norm limits and pass the double-centralizer identities to
these limits.

\textbf{Theorem 2.1 (Kasparov's technical theorem).} Let \(J\) be a
\(\sigma\)-unital C*-algebra. Let \(A_1,A_2\subset M(J)\) be
\(\sigma\)-unital C*-subalgebras, and let \(\Delta\subset M(J)\) be a
separable linear subspace. Assume \[
A_1A_2\subset J,\qquad [\Delta,A_1]\subset A_1.
\tag{2.1}
\] There are \(M,N\in M(J)\) satisfying \[
\begin{gathered}
0\leq M,N\leq1,\qquad M+N=1,\\
MA_1\subset J,\qquad NA_2\subset J,\\
[M,\Delta]\subset J.
\end{gathered}
\tag{2.2}
\]

The identity belongs to the multiplier algebra; the theorem does not
require either \(A_i\) to contain it.

\textbf{Proof for separable data.} Enlarge \(\Delta\) by its adjoints.
This still derives \(A_1\), since \([d^*,a]=-[d,a^*]^*\). Choose
norm-compact sets \(X_1,X_2,Y\), of elements of norm at most one, whose
linear spans are dense in \(A_1,A_2,\Delta\). Such a set is obtained
from a bounded dense sequence by multiplying its \(j\)-th term by
\(1/j\) and adjoining zero. Include adjoints when needed.

Put \(\varepsilon_n=2^{-n}\). By Theorem 1.1 choose positive
contractions \(u_n\in A_1\) such that \[
\|u_nx-x\|<\varepsilon_n\quad(x\in X_1),\qquad
\|[u_n,y]\|<\varepsilon_n\quad(y\in Y).
\tag{2.3}
\] They may be chosen as an approximate identity of \(A_1\).

Using Lemmas 1.2--1.3 on \(J\), choose an increasing approximate
identity \(v_n\in J\), with \(v_0=0\), and set \[
b_n=(v_n-v_{n-1})^{1/2}.
\tag{2.4}
\] Require \[
\begin{aligned}
\|(1-v_n)u_jx\|&<\varepsilon_n^4
&& (x\in X_2,\ 1\leq j\leq n+1),\\
\|[b_n,w]\|&<\varepsilon_n
&& (w\in X_1\cup X_2\cup Y).
\end{aligned}
\tag{2.5}
\] The first requirement tests compact subsets of \(J\), because
\(u_jX_2\subset A_1A_2\subset J\). The second can be arranged
prospectively: if \(\eta_n\) is the square-root tolerance in Lemma 1.3
for \(\varepsilon_n\), require the commutators of \(v_n\) to be smaller
than half both \(\eta_n\) and \(\eta_{n+1}\). Then the commutators of
\(v_n-v_{n-1}\) are smaller than \(\eta_n\). The multiplier tests derive
the ideal \(J\), so Lemma 1.2 applies. Include an increasing dense
family of \(J\)-tests to keep \(v_n\) an approximate identity.

Define partial sums \[
N_r=\sum_{n=1}^r b_nu_nb_n.
\tag{2.6}
\] Each summand is positive and belongs to \(J\). Since \(u_n\leq1\), \[
0\leq N_s-N_r\leq v_s-v_r\leq1-v_r\qquad(s\geq r),
\tag{2.7}
\] and \(N_r\leq v_r\leq1\). For \(z\in J\), \[
\|(N_s-N_r)z\|^2
\leq\|z^*(N_s-N_r)z\|
\leq\|z^*(1-v_r)z\|\longrightarrow0.
\tag{2.8}
\] Applying the same argument to \(z^*\) gives convergence on the other
side. Thus (2.6) converges strictly to a positive contraction
\(N\in M(J)\). Put \(M=1-N\).

We now prove all three conclusions in (2.2), with norm estimates. For
\(x\in X_2\) and \(n\geq2\), both \(v_n\) and \(v_{n-1}\) approximate
\(u_nx\) by the first line of (2.5); this is why that line included
\(j=n+1\). Thus \[
\begin{aligned}
\|(v_n-v_{n-1})u_nx\|&<17\varepsilon_n^4,\\
\|b_nu_nx\|^2
&=\|(u_nx)^*(v_n-v_{n-1})(u_nx)\|
<17\varepsilon_n^4.
\end{aligned}
\tag{2.9}
\] Commuting the last \(b_n\) past \(x\), we obtain \[
\|b_nu_nb_nx\|
\leq\sqrt{17}\,\varepsilon_n^2+\varepsilon_n.
\tag{2.10}
\] The series \(Nx=\sum b_nu_nb_nx\) is therefore norm-convergent, apart
from one harmless first term, and all its terms lie in \(J\). Hence
\(NA_2\subset J\), by density and boundedness.

For \(x\in X_1\), expand the complementary summand: \[
\begin{gathered}
(b_n^2-b_nu_nb_n)x\\
 =b_n(x-u_nx)b_n\\
 \quad+b_n[b_n,x]\\
 \quad-b_nu_n[b_n,x].
\end{gathered}
\tag{2.11}
\] Its norm is at most \(3\varepsilon_n\), by (2.3) and (2.5). The
norm-convergent sum of these elements belongs to \(J\). Strictly, that
sum is \[
\left(\sum b_n^2-\sum b_nu_nb_n\right)x=(1-N)x.
\] Thus \(MA_1\subset J\).

Finally for \(y\in Y\), \[
[b_nu_nb_n,y]
=b_nu_n[b_n,y]+b_n[u_n,y]b_n+[b_n,y]u_nb_n,
\tag{2.12}
\] of norm at most \(3\varepsilon_n\). Its norm-convergent sum lies in
\(J\) and equals \([N,y]\) by strict continuity of multiplication by a
fixed multiplier. Density gives \([N,\Delta]\subset J\), and
\([M,\Delta]=-[N,\Delta]\). This proves the separable case. \(\square\)

\textbf{Reduction to the stated generality.} Choose strictly positive
\(h_0\in J\), \(h_1\in A_1\), \(h_2\in A_2\). Inside \(A_1\), begin with
\(C^*(h_1)\) and repeatedly adjoin commutators with a countable dense
family of \(\Delta+\Delta^*\). The closed union is a separable
subalgebra \(B_1\) containing \(h_1\) and derived by \(\Delta\). Put
\(B_2=C^*(h_2)\).

Inside \(J\), start with the separable algebra generated by \(h_0\) and
\(B_1B_2\). Repeatedly adjoin its left and right products by
\(B_1,B_2,\Delta,\Delta^*\). Their closed union \(J_0\) is separable and
is preserved on both sides by all those multipliers. It contains the
strictly positive element \(h_0\) of the original \(J\).

There is a canonical embedding \(M(J_0)\subset M(J)\). Here are the
details that make the reduction legitimate. An approximate identity of
\(J_0\) is also one for \(J\): it approximates \(h_0\), and the density
of \(h_0J\) and \(Jh_0\) extends that approximation to every element of
\(J\). A bounded strict Cauchy family on \(J_0\) is then strict Cauchy
on \(J\), by first multiplying a \(J\)-test by a fixed
approximate-identity element of \(J_0\). Consequently the bounded strict
approximation \(me_\lambda\in J_0\) of a multiplier \(m\in M(J_0)\)
extends to a multiplier of \(J\). Its multiplication and adjoint are
preserved by strict limits. Restriction to \(J_0\) proves injectivity
and identifies every original multiplier preserving \(J_0\) with its
extension.

Apply the proved separable case to \(J_0,B_1,B_2,\Delta\). Its \(M,N\)
lie in \(M(J)\), with \(MB_1,NB_2,[M,\Delta]\subset J_0\). If
\(e_n=f_n(h_1)\) is an approximate identity of \(A_1\), then
\(e_n\in B_1\), and for \(a\in A_1\), \[
Ma=\lim_n Me_na\in J.
\tag{2.13}
\] Likewise \(NA_2\subset J\), using \(h_2\). The commutator conclusion
already lies in \(J_0\subset J\). This completes the theorem.
\(\square\)

\subsection{3. Grading and square roots}\label{grading-and-square-roots}

For graded \(J\), its grading extends to \(M(J)\). A graded subspace
derives a graded algebra when the graded commutator is in that algebra
for homogeneous inputs.

\textbf{Theorem 3.1 (graded technical theorem).} In Theorem 2.1, suppose
\(J,A_1,A_2,\Delta\) are graded and the derivating condition uses graded
commutators. Then \(M,N\) can be chosen even. All conclusions (2.2)
hold, with graded commutators.

\textbf{Proof.} A graded \(\sigma\)-unital algebra has an even strictly
positive element: replace \(h\) by \(h+\gamma(h)\). Here is the
domination argument. If \(k\ge h\ge0\), set
\(e_\varepsilon=k(k+\varepsilon)^{-1}\), in the unitization when
necessary. Functional calculus and order give \[
0\le(1-e_\varepsilon)h(1-e_\varepsilon)
\le(1-e_\varepsilon)k(1-e_\varepsilon)
\le\frac{\varepsilon}{4}1.
\] The last bound is
\(\sup_{t\ge0}t\varepsilon^2/(t+\varepsilon)^2=\varepsilon/4\). Hence
\(\|(1-e_\varepsilon)h\|\to0\), by applying the C*-identity to
\(h^{1/2}(1-e_\varepsilon)\) and then multiplying by \(h^{1/2}\).
Density of \(hA\) gives \(e_\varepsilon a\to a\) for every \(a\); taking
adjoints gives the right convergence. Since \(e_\varepsilon a\in kA\),
that range is dense. Apply this to \(k=h+\gamma(h)\). Its
functional-calculus approximate identity is even.

Use the bounded-functional vector representations and Hilbert support
projections constructed in Theorem 1.1, with an even approximate
identity. For homogeneous \(d\) and \(a\), the identity
\(\pi(d)\pi(a)\eta=(-1)^{|d||a|}\pi(a)\pi(d)\eta+\pi([d,a]_{\mathrm{gr}})\eta\)
shows that \(\overline{\pi(A)H}\) is invariant under \(\pi(d)\); the
same holds for \(d^*\). Thus this subspace reduces \(\pi(d)\), and its
projection commutes with it. Represented approximate identities converge
strongly to that projection. Their commutators with \(d\) therefore
converge strongly to zero in every such representation;
vector-coefficient testing gives Banach weak convergence. These are also
the graded commutators, since the approximate identity is even. The
convex-separation argument of Theorem 1.1 now gives even convex
combinations, and the increasing construction and square-root functional
calculus retain that parity.

In the separable reduction use grading-invariant subalgebras and
homogeneous dense families, adjoining graded commutators. This keeps
\(B_1\) invariant and graded-derived. Choose the \(u_n,v_n,b_n\) even.
Every commutator involving one of them is then ordinary, so each
estimate (2.3)--(2.12) applies unchanged. The strict sum \(N\) and its
complement \(M\) are even. Their commutators with homogeneous tests are
the required graded commutators. \(\square\)

\textbf{Corollary 3.2 (the square-root partition).} Under the hypotheses
of Theorem 3.1, \[
\begin{gathered}
M^{1/2}A_1\subset J,\qquad N^{1/2}A_2\subset J,\\
[M^{1/2},\Delta]\subset J,\qquad
[N^{1/2},\Delta]\subset J.
\end{gathered}
\tag{3.1}
\]

\textbf{Proof.} For \(a\in A_1\), \((M^{1/2}a)^*(M^{1/2}a)=a^*Ma\in J\).
Passing to \(M(J)/J\) shows that \(M^{1/2}a\) has zero quotient norm, so
it lies in \(J\). The same argument applies to \(N\) and \(A_2\). In the
quotient, the images of \(M,N\) commute with \(\Delta\). Polynomial
approximation to the square root preserves that commutation. Lifting
back gives the two commutator conclusions. \(\square\)

This is the form used to combine product operators:
\(M^{1/2}F_1+N^{1/2}F_2\). The partition disposes of defects; the next
lesson must also verify the connection and positivity conditions for
that product.

\subsection{4. A multiplier lifting bridge and the simpler
separation}\label{a-multiplier-lifting-bridge-and-the-simpler-separation}

The absence of commutator tests admits another proof. To make that proof
usable, we supply the needed lifting statement here.

\textbf{Theorem 4.1 (multiplier quotient lifting).} If \(D\) is
\(\sigma\)-unital and \(I\) is an ideal, the quotient map extends to a
surjective homomorphism \[
\widetilde q:M(D)\longrightarrow M(D/I).
\tag{4.1}
\] Every positive contraction in the target has a positive contractive
lift.

\textbf{Proof.} Write \(R=D/I\). Every multiplier \(m_0\) of \(D\)
preserves \(I\). Indeed, for \(i\in I\) and an approximate identity
\(a_\lambda\) of \(I\), the elements \((m_0i)a_\lambda\) belong to
\(I\), since \(m_0i\in D\), and they converge to \(m_0i\), since
\(ia_\lambda\to i\). Taking adjoints gives the other side. Thus the two
multiplier actions descend to \(R\). This defines \(\widetilde q\); on
bounded families it is strictly continuous, because its products with
\(q(d)\) are the quotients of the products with \(d\).

Take \(0\leq m\leq1\) in \(M(R)\). Begin with the nested approximate
identity (1.3) in \(D\). Its image is an approximate identity of \(R\).
The multiplier \(m\) derives \(R\). By Lemma 1.2 choose convex
combinations \(v_n\in D\), increasing as in (1.4), such that, for
\(b_n=(v_n-v_{n-1})^{1/2}\), \[
\|[q(b_n),m]\|<2^{-n}.
\tag{4.2}
\] Use Lemma 1.3 to prescribe the commutator tolerances for both
consecutive \(v_n\), exactly as in (2.5).

For each \(n\), an approximate identity \(w_k\) of \(R\) satisfies \[
\|(w_k^{1/2}mw_k^{1/2}-m)q(b_n)\|\longrightarrow0.
\] This follows by approximating \(q(b_n)\) and \(mq(b_n)\) by the same
approximate identity; square roots are themselves an approximate
identity. Choose \(k=k(n)\) making the norm less than \(2^{-n}\). Lift
the positive contraction \(w_k^{1/2}mw_k^{1/2}\in R\) to a positive
contraction \(a_n\in D\). Such an element lift follows by first lifting
self-adjointly and then clipping by the continuous function
\(t\mapsto\min(1,\max(0,t))\).

The positive sum \[
L=\sum_{n=1}^\infty b_na_nb_n
\tag{4.3}
\] converges strictly to a positive contraction in \(M(D)\): its partial
sums and tails are dominated as in (2.7), and (2.8) proves strict
convergence. By strict continuity,
\(\widetilde q(L)=\sum q(b_n)q(a_n)q(b_n)\). The differences \[
\begin{aligned}
q(b_n)q(a_n)q(b_n)-q(b_n)mq(b_n),\\
q(b_n)mq(b_n)-q(b_n)^2m
\end{aligned}
\tag{4.4}
\] belong to \(R\) and each have norm less than \(2^{-n}\), by the
selection of \(a_n\) and (4.2). Their sums consequently converge in norm
in \(R\). Meanwhile \(\sum q(b_n)^2m=m\) strictly. Hence
\(\widetilde q(L)-m\in R\).

Choose a self-adjoint \(d\in D\) with \(q(d)=m-\widetilde q(L)\). Then
\(L+d\) is a self-adjoint multiplier lifting \(m\). Clipping it to
\([0,1]\) gives a positive contractive multiplier lift, since functional
calculus commutes with the quotient and fixes \(m\). Finally every
multiplier is a linear combination of positive multipliers, by splitting
its real and imaginary parts into positive and negative parts. This
proves surjectivity. \(\square\)

\textbf{Corollary 4.2 (separation without commutator tests).} Theorem
2.1 with \(\Delta=0\) follows from Theorem 4.1.

\textbf{Proof.} Let \(D=C^*(J,A_1,A_2)\subset M(J)\). This algebra is
\(\sigma\)-unital. To verify that point, add strictly positive elements
of \(J,A_1,A_2\). The sum is positive and its functional-calculus
approximate identity approximates each generator algebra; for example
domination of the strictly positive element \(h_i\) gives
\(\|(1-e_n)h_i^{1/2}\|\to0\), then density of \(h_iA_i\) gives
approximation on \(A_i\). Approximation extends to their generated
algebra.

The quotient is \(q(D)=q(A_1)\oplus q(A_2)\), since these two
subalgebras are orthogonal. In its multiplier algebra choose the central
projection \(r=(0,1)\). Theorem 4.1 lifts \(r\) to a positive
contraction \(M\in M(D)\). Restricting multipliers to the essential
ideal \(J\) embeds \(M(D)\) in \(M(J)\): multiplication on \(J\)
determines a multiplier, and essentiality detects its kernel. Put
\(N=1-M\). Their quotient products with \(q(A_1)\) and \(q(A_2)\) vanish
in the required order, so \(MA_1,NA_2\subset J\). \(\square\)

This proof has supplied the multiplier lifting, rather than presupposing
it. It does not produce a multiplier commuting modulo \(J\) with a
specified \(\Delta\); that additional constraint is the content of the
series selection in Section 2.

\subsection{5. Extending an operator homotopy from an
ideal}\label{extending-an-operator-homotopy-from-an-ideal}

Let \(A\) be separable and graded, \(I\subset A\) a graded ideal, and
\(E\) a countably generated graded Hilbert \(B\)-module. Let
\(\phi:A\to\mathcal L(E)\) be graded. Write \(K=\mathcal K(E)\) and
define the two pseudolocal algebras \[
\begin{aligned}
C&=\{x:[x,\phi(a)]_{\mathrm{gr}}\in K\ (a\in A)\},\\
D&=\{x:[x,\phi(i)]_{\mathrm{gr}}\in K\ (i\in I)\}.
\end{aligned}
\tag{5.1}
\] Use homogeneous components in this notation. Their locally compact
ideals are \[
\begin{aligned}
J_A&=\{x\in C:x\phi(a),\phi(a)x\in K\ (a\in A)\},\\
J_I&=\{x\in D:x\phi(i),\phi(i)x\in K\ (i\in I)\}.
\end{aligned}
\tag{5.2}
\] These are closed graded C*-ideals in the respective algebras, by the
commutator Leibniz rule and adjoints, as in the quotient formulation of
the preceding cycle lesson.

\textbf{Lemma 5.1 (pseudolocal decomposition).} One has \(D=C+J_I\).
Consequently inclusion induces a surjective map \[
C/J_A\longrightarrow D/J_I.
\tag{5.3}
\]

\textbf{Proof.} Take homogeneous \(x\in D\). For homogeneous
\(a\in A,i\in I\), the graded Leibniz rule gives \[
[x,\phi(a)]_{\mathrm{gr}}\phi(i)
=[x,\phi(ai)]_{\mathrm{gr}}
-(-1)^{|x||a|}\phi(a)[x,\phi(i)]_{\mathrm{gr}}\in K.
\tag{5.4}
\] Adjoints give compactness in the other order. Thus the algebra \[
A_1=C^*(K,\phi(I))
\] has compact products with \[
A_2=C^*(K,\{[x,\phi(a)]_{\mathrm{gr}}:a\in A\}).
\tag{5.5}
\] Indeed the localized commutators, their adjoints, and their finite
products annihilate \(\phi(I)\) modulo \(K\), by (5.4), and products
involving \(K\) are compact.

Both \(A_1,A_2\) are \(\sigma\)-unital: \(K\) is \(\sigma\)-unital, and
each adjoins a separable set of generators. Explicitly, if \(h_K\) is
strictly positive in \(K\) and \(s_j\) is a bounded countable generating
set for the adjoined part, the positive sum \[
h_K+\sum_j2^{-j}(s_j^*s_j+s_js_j^*)
\tag{5.6}
\] has functional-calculus approximate identity converging on \(K\) and
each \(s_j,s_j^*\), by domination. It therefore approximates the
generated algebra. Take homogeneous generators to make the sum even.

The separable graded subspace spanned by \(x,x^*,\phi(A)\) derives
\(A_1\). Commutators with \(K\) remain in \(K\); commutators of
\(\phi(A)\) with \(\phi(I)\) remain in \(\phi(I)\); and the commutators
of \(x,x^*\) with \(\phi(I)\) are compact. Apply Theorem 3.1 in
\(M(K)=\mathcal L(E)\). It gives even \(M,N\) such that \[
M\phi(I)\subset K,\quad
N[x,\phi(A)]_{\mathrm{gr}}\subset K,\quad
[M,x],[M,\phi(A)]\subset K.
\tag{5.7}
\] The graded product rule shows \(Nx\in C\). The other part \(Mx\)
belongs to \(J_I\): commuting \(x\) across \(\phi(i)\) leaves a compact
error and a factor \(M\phi(i)\), and the same argument with adjoints
gives the other side. Its pseudolocality for \(I\) follows from the same
product rule. Thus \(x=Nx+Mx\in C+J_I\).

Decompose an arbitrary \(x\) into its homogeneous parts. This proves the
equality and the surjectivity in (5.3); the inclusion maps \(J_A\) into
\(J_I\). \(\square\)

The localized compactness conditions in (5.2) must be stated inside the
pseudolocal algebras. A one-sided compactness set of arbitrary
adjointable operators need not itself be a C*-algebra.

\textbf{Lemma 5.2 (lifting an involution path).} For a surjective graded
homomorphism between unital C*-algebras, a norm-continuous path of odd
self-adjoint unitaries lifts to such a path, starting at any prescribed
odd self-adjoint unitary lift.

\textbf{Proof.} First recall an elementary Banach quotient argument.
Approximate a continuous quotient path by finite polygonal paths,
retaining its endpoints and with geometrically decreasing uniform
errors. Lift the differences of consecutive polygonal paths by lifting
their finitely many vertices to elements of at most twice their quotient
norm and interpolating. The differences vanish at the endpoints, so
their lifts can vanish there too. The uniform sum of those correction
paths converges, giving a lift with prescribed endpoint values.

If instead only the initial value zero is prescribed, and the quotient
path has norm at most \(\delta\), begin with a polygonal approximation
of error at most \(\delta/4\). Its vertex lifts have norms at most
\(2\delta\). Subsequent errors at most \(\delta/2^{n+2}\) yield
uniformly summable correction lifts of total norm at most \(2\delta\). A
bound \(8\delta\) therefore holds for the full lift. Contractive
averaging with the adjoint and grading makes every lift self-adjoint and
odd when the quotient path has those properties.

Subdivide the involution path until each piece stays within a
sufficiently small \(\delta\) of its initial value. Lift its difference
from that initial value by this argument, starting at zero.

Choose the subdivision so these lifted differences have norm less than
\(1/2\). Adding one to the current self-adjoint unitary lift gives an
odd self-adjoint invertible path \(h_t\). Replace it by \[
h_t|h_t|^{-1}.
\tag{5.8}
\] It is an odd self-adjoint unitary path. Its quotient is the original
involution path, because (5.8) fixes a self-adjoint unitary. At the
start it equals the prescribed lift. Continue from the new endpoint on
the next subinterval; the resulting finitely many paths glue.
\(\square\)

For logical independence, the elementary polygonal lifting used here
needs no result about KK. Its proof is the summable finite-vertex
correction argument, valid in every Banach quotient; the self-adjoint
and odd parts are obtained by contractive averaging.

\textbf{Theorem 5.3 (ideal homotopy extension).} Let \((E,\phi,F)\) be a
cycle over \((A,B)\), with \(A\) separable. Suppose \(F_t\) is a
norm-continuous operator homotopy for its restriction to \(I\), starting
with \(F_0=F\). There is an operator homotopy \(G_t\) for the original
representation of \(A\), with \(G_0=F\), and \[
(G_t-F_t)\phi(i),\quad\phi(i)(G_t-F_t)\in\mathcal K(E)
\qquad(t\in[0,1],\ i\in I).
\tag{5.9}
\]

\textbf{Proof.} The quotient formulation of cycle conditions identifies
the images of \(F_t\) in \(D/J_I\) as odd self-adjoint unitaries. The
image of \(F\) in \(C/J_A\) is one such lift of their initial value. By
Lemmas 5.1--5.2, lift the whole path to odd self-adjoint unitaries in
\(C/J_A\). The Banach-space path lifting argument then lifts this to a
norm-continuous odd path \(G_t\in C\), starting at the original \(F\).
Self-adjointness is required only modulo \(J_A\), so the prescribed
\(F\) need not be exactly self-adjoint.

In \(C/J_A\), \(G_t^2=1\) and \(G_t=G_t^*\), so its square and adjoint
defects are locally compact for all \(a\in A\). Membership in \(C\)
gives its compact graded commutators. It is therefore the required
operator homotopy. Its image in \(D/J_I\) equals that of \(F_t\),
proving (5.9). \(\square\)

\subsection{6. Examples and exercises}\label{examples-and-exercises}

If \(A_2=0\), take \(M=0,N=1\). This meets every condition, including
all commutators. If \(A_1\subset J\), take \(M=1,N=0\): then
\(MA_1\subset J\), \(NA_2=0\), and all commutators vanish. The
interesting case is when the two algebras are only orthogonal after
passing to the quotient.

For a concrete compact-operator example, take
\(J=\mathcal K(L^2(S^1))\), the Hardy projection \(P\), and
multiplication representation \(\pi\) of \(C(S^1)\). Put \[
A_1=\mathcal K(L^2(S^1))+\mathbb C(1-P),\quad
A_2=\mathbb CP,\quad \Delta=\pi(C(S^1)).
\tag{6.1}
\] These are separable and \(\sigma\)-unital. Their products are
compact. The compact Hardy commutators show that \(\Delta\) derives
\(A_1\). Here the explicit partition \(M=P,N=1-P\) already satisfies the
theorem. The general construction replaces this visible projection when
no exact geometric projection is available.

\textbf{6.1. Finite-rank quasicentrality.} For a fixed
\(T\in\mathcal B(H)\), where \(H\) is separable, construct positive
finite-rank contractions \(u_n\), each a convex combination of
finite-rank projections from an increasing sequence converging strongly
to \(1\), such that \(\|[u_n,T]\|\to0\).

\textbf{6.2. Why projections are too restrictive.} Let \(S\) be the
unilateral shift. Prove that finite-rank projections \(P_n\to1\)
strongly cannot satisfy \(\|[P_n,S]\|\to0\). Increasingness is not
needed for this obstruction.

\textbf{6.3. Separation through multiplier lifting.} For \(\Delta=0\),
form \(D=C^*(J,A_1,A_2)\) and use the central projection in \(M(D/J)\)
to deduce the technical theorem. Verify the \(\sigma\)-unitality and
essential-ideal steps.

\textbf{6.4. Even partitions.} Prove the graded theorem, explaining why
merely averaging an arbitrary ungraded partition is insufficient when
the given derivating hypothesis is only graded.

\subsection{7. Solutions}\label{solutions}

\textbf{Solution to 6.1.} Let \(p_n\) project onto the first \(n\)
vectors of an orthonormal basis. They form an approximate identity of
\(\mathcal K(H)\). The subspace spanned by \(T,T^*\) derives that ideal,
since multiplication by a bounded operator preserves compactness. Apply
Theorem 1.1 to successive tails, making the two commutators smaller than
\(1/n\). The selected convex combinations are positive finite-rank
contractions. Choosing each tail beyond \(n\) retains strong convergence
to \(1\), since all projections in the combination fix the first \(n\)
basis vectors. In particular the commutator with \(T\) tends to zero in
norm.

\textbf{Solution to 6.2.} Let \(p_0\) project onto the first basis
vector \(e_0\), so \(SS^*=1-p_0\). For one finite-rank projection \(P\),
put \(C=PSP\) on the finite-dimensional space \(PH\), and
\(\delta=\|[P,S]\|\). Its two off-diagonal blocks have norms at most
\(\delta\). Since \(S^*S=1\), \[
0\leq P-C^*C=PS^*(1-P)SP,\qquad
\|P-C^*C\|\leq\delta^2.
\tag{7.1}
\] If \(\delta<1\), \(C\) is invertible on \(PH\). Its polar
decomposition gives \(\|P-CC^*\|=\|P-C^*C\|\leq\delta^2\). But \[
P-CC^*
=Pp_0P+PS(1-P)S^*P\geq Pp_0P.
\tag{7.2}
\] Hence \(\|Pe_0\|^2=\|Pp_0P\|\leq\delta^2\). Strong convergence
\(P_ne_0\to e_0\) contradicts \(\delta_n\to0\). This also explains why
convex positive contractions, rather than projections, appear in the
technical theorem.

\textbf{Solution to 6.3.} The positive sum of strictly positive elements
of \(J,A_1,A_2\) is strictly positive in \(D\), by the domination and
generator argument in Corollary 4.2. The quotient \(D/J\) is the
orthogonal sum \(q(A_1)\oplus q(A_2)\). Lift its central multiplier
projection \((0,1)\) by Theorem 4.1 to a positive contraction
\(M\in M(D)\). The ideal \(J\) is essential in \(D\), since
\(D\subset M(J)\) and a multiplier annihilating \(J\) is zero. Thus
restriction embeds \(M(D)\) in \(M(J)\). Quotient multiplication gives
\(MA_1\subset J\) and \((1-M)A_2\subset J\). Taking \(N=1-M\) proves
every conclusion; there are no commutator tests.

\textbf{Solution to 6.4.} Graded derivation does not imply ordinary
derivation on two odd inputs, so one cannot simply invoke the ungraded
theorem and average its answer. Instead begin with even strictly
positive elements and even functional-calculus approximate identities.
Their commutators with homogeneous tests are ordinary commutators and
graded commutators simultaneously. In each vector-coefficient
representation used in Theorem 1.1, the represented approximate identity
tends strongly to the projection onto the closure of the represented
algebra acting on the Hilbert space. For a homogeneous test and a
homogeneous algebra element, the graded derivation identity proves
invariance of this subspace under the test and its adjoint, just as in
that proof; the sign changes only the first term. Thus the subspace
reduces the test and its support projection commutes with it. Since the
approximate identity is even, its ordinary and graded commutators
coincide. Testing vector coefficients proves the required Banach weak
convergence and hence the convexity step. Retain evenness through the
convex combinations, the separable invariant reduction and the square
roots \(b_n\). All factors used in the estimates are even, so no extra
signs occur. The resulting positive strict sum and its complement are
even and satisfy the graded conditions. This verifies each step of
Theorem 3.1.

\subsection{What this lesson does not
prove}\label{what-this-lesson-does-not-prove}

Continuous functional calculus, the GNS construction and the
compact-module multiplier identification are the exact earlier
prerequisites stated at the beginning. The bounded-functional vector
representation and the support-projection argument were proved within
Theorem 1.1. They provide the framework; all convex quasicentrality,
separation, quotient multiplier lifting and ideal-homotopy extension
arguments were proved here.

The construction and uniqueness of the Kasparov product, its
associativity, boundary products and the equivariant technical theorem
require further proofs in the following lessons. The partition theorem
alone does not assert them.

\subsection{References}\label{references}

\begin{itemize}
\item
  B. Blackadar,
  \href{https://www.bruceblackadar.com/Mathematics/book6.pdf}{\emph{K-Theory
  for Operator Algebras}, author-posted corrected second edition}. The
  \href{https://www.bruceblackadar.com/mathpubs.html}{author identifies
  this freely downloadable version} as a slightly corrected second
  edition. The source retains its author copyright and the usage terms
  stated on that page. Sections 12.4 and 14.6 supply the free source
  comparison for quasicentrality and the strict-series partition. The
  convex separation, increasing choices, graded construction and
  multiplier lifting used here are proved in Sections 1--4 above; the
  proof uses the preceding results identified above.
\item
  \href{KT-KK-05.html\#lemma-2-0-compact-module-operators-and-their-multipliers}{Lesson
  05, Lemma 2.0}, for \(\sigma\)-unitality of \(\mathcal K(E)\),
  \(\mathcal L(E)=M(\mathcal K(E))\), and bounded strict limits.
\item
  \emph{Foundations of von Neumann algebras},
  \href{https://kokunoyumeto.github.io/open-math-courses/courses/foundations-of-von-neumann-algebras/c-star-algebras-continuous-functional-calculus-automatic-continuity-positive-cones.html}{Continuous
  functional calculus}, Theorems 5.1 and 6.1.
\item
  B. Blackadar,
  \href{https://www.bruceblackadar.com/Mathematics/Cycr.pdf}{\emph{Operator
  Algebras}, freely readable corrected author edition dated February 8,
  2017}, Sections II.6.3--II.6.4, for positive state extension and GNS
  comparison. Theorem 1.1 gives the complete bounded-functional
  construction and weak-support proof used here. The proof does not
  require Jordan decomposition or a universal-bidual theorem. The author
  retains copyright and the terms stated on the linked author
  publications page.
\end{itemize}

\end{document}
